%% =====================================================================
\chapter{Prescriptions, Vital Signs and IoT Devices}\label{ch:rxvitals}

\section{Prescriptions}

\textbf{Files.} \file{src\codesl{}models\codesl{}Prescription\codedot{}js},
\file{src\codesl{}modules\codesl{}prescriptions\codesl{}prescriptions\codedot{}routes\codedot{}js}.

A doctor with the patient's consent prescribes a medication with dosage,
frequency, instructions, dates and a number of allowed refills. The patient sees
active and past medications and can request a refill; the doctor approves or
denies it.

\begin{lstlisting}[language=JavaScript,caption={Refill request rules (\file{prescriptions.routes.js})}]
router.post('/:id/refill-request', requireRole('patient'), validate(idParam), asyncHandler(async (req, res) => {
  const rx = await Prescription.findOne({ _id: req.params.id, patientId: req.user.userId });
  if (!rx) throw notFound('Prescription not found');
  if (rx.status !== 'active') throw conflict('Only active prescriptions can be refilled', 'INVALID_TRANSITION');
  if (rx.refillRequest.status === 'pending') throw conflict('A refill request is already pending', 'ALREADY_PENDING');
  if (rx.refillsUsed >= rx.refillsAllowed) throw conflict('No refills remain on this prescription. Book an appointment for a new one.', 'NO_REFILLS');

  rx.refillRequest = { status: 'pending', requestedAt: new Date(), decidedAt: null };
  await chain(rx, 'REFILL_REQUESTED', req);
  await notificationService.notify(rx.doctorId, {
    type: 'warning', title: 'Refill request', message: `A patient requested a refill of ${rx.medication} ${rx.dosage}.`, page: 'prescriptions',
  });
  res.json({ success: true, message: 'Refill requested', data: rx });
}));
\end{lstlisting}

\begin{lstlisting}[language=JavaScript,caption={The doctor decides (\file{prescriptions.routes.js})}]
    if (action === 'approve_refill' || action === 'deny_refill') {
      if (rx.refillRequest.status !== 'pending') throw conflict('There is no pending refill request', 'INVALID_TRANSITION');
      const approved = action === 'approve_refill';
      rx.refillRequest.status = approved ? 'approved' : 'denied';
      rx.refillRequest.decidedAt = new Date();
      if (approved) rx.refillsUsed += 1;
      await chain(rx, approved ? 'REFILL_APPROVED' : 'REFILL_DENIED', req);
      // ... notify the patient
    } else {
      if (rx.status !== 'active') throw conflict(`The prescription is already ${rx.status}`, 'INVALID_TRANSITION');
      rx.status = action === 'complete' ? 'completed' : 'cancelled';
      await chain(rx, action === 'complete' ? 'COMPLETE' : 'CANCEL', req);
    }
\end{lstlisting}

The doctor can only act on prescriptions they wrote (the query includes
\code{doctorId}). All rule violations return 409 with a precise message, which
the dashboard shows as-is. Every change is chained in the ledger, so the
dispensing history of a controlled drug can be audited.

\section{Vital signs}

\textbf{Files.} \file{src\codesl{}models\codesl{}VitalReading\codedot{}js}, \file{src\codesl{}models\codesl{}Device\codedot{}js},
\file{src\codesl{}modules\codesl{}vitals\codesl{}vitals\codedot{}routes\codedot{}js}.

\begin{lstlisting}[language=JavaScript,caption={Normal ranges and the abnormal check (\file{src/models/VitalReading.js})}]
const VITAL_RANGES = {
  heartRate: { min: 50, max: 110, unit: 'bpm', label: 'Heart rate' },
  systolic: { min: 90, max: 140, unit: 'mmHg', label: 'Systolic BP' },
  diastolic: { min: 60, max: 90, unit: 'mmHg', label: 'Diastolic BP' },
  spo2: { min: 94, max: 100, unit: '%', label: 'Blood oxygen' },
  temperature: { min: 36.1, max: 37.8, unit: 'C', label: 'Temperature' },
  glucose: { min: 70, max: 180, unit: 'mg/dL', label: 'Blood glucose' },
};

/** Returns the names of the fields in `reading` that fall outside VITAL_RANGES. */
function findAbnormal(reading) {
  return Object.entries(VITAL_RANGES)
    .filter(([field, range]) => reading[field] !== null && reading[field] !== undefined
      && (reading[field] < range.min || reading[field] > range.max))
    .map(([field]) => field);
}
\end{lstlisting}

(The temperature unit in the source is the degree sign, \textdegree C.) The
ranges are sent to the frontend with every vitals response, so the charts and
tiles colour values with the same thresholds the server uses for alerts.

\begin{lstlisting}[language=JavaScript,caption={Saving a reading and raising alerts (\file{vitals.routes.js})}]
async function saveReading(patientId, values, source, deviceId = null) {
  const reading = await VitalReading.create({
    patientId, deviceId, source, ...values, abnormal: findAbnormal(values), recordedAt: values.recordedAt || new Date(),
  });

  if (reading.abnormal.length) {
    const details = reading.abnormal.map((f) => `${VITAL_RANGES[f].label} ${reading[f]} ${VITAL_RANGES[f].unit}`).join(', ');
    const patient = await User.findById(patientId).select('fullName');
    const doctors = await AccessGrant.find({ patientId, status: 'active', expiresAt: { $gt: new Date() } }).distinct('doctorId');
    await notificationService.notify(patientId, {
      type: 'danger', title: 'Abnormal vital sign', message: `Out of normal range: ${details}.`, page: 'metrics',
    });
    await notificationService.notify(doctors, {
      type: 'danger', title: `Vitals alert: ${patient.fullName}`, message: `Out of normal range: ${details}.`, page: 'iot',
    });
  }
  return reading;
}
\end{lstlisting}

Alerts go to the patient and to every doctor who \emph{currently} has consent:
the consent model decides who is told, not just who can look.

\section{IoT devices}

\subsection{Registration and keys}

An admin registers a device (serial number, type, patient). The server generates
a random API key, stores only its SHA-256, and shows the key once:

\begin{lstlisting}[language=JavaScript,caption={Device model (\file{src/models/Device.js})}]
const deviceSchema = new mongoose.Schema({
  serialNumber: { type: String, required: true, unique: true, trim: true },
  name: { type: String, required: true, trim: true },
  type: { type: String, enum: ['heart_monitor', 'bp_monitor', 'pulse_oximeter', 'thermometer', 'multi_parameter'], required: true },
  patientId: { type: mongoose.Schema.Types.ObjectId, ref: 'User', required: true, index: true },
  apiKeyHash: { type: String, required: true, unique: true },
  status: { type: String, enum: ['active', 'inactive'], default: 'active' },
  lastSeenAt: { type: Date, default: null },
  registeredBy: { type: mongoose.Schema.Types.ObjectId, ref: 'User', required: true },
}, { timestamps: true });
\end{lstlisting}

\subsection{Ingestion}

\begin{lstlisting}[language=JavaScript,caption={POST /api/devices/readings (\file{vitals.routes.js})}]
router.post('/devices/readings', validate({ body: readingSchema }), asyncHandler(async (req, res) => {
  const key = req.get('x-device-key');
  if (!key) throw unauthorized('Device key required', 'NO_DEVICE_KEY');
  const device = await Device.findOne({ apiKeyHash: crypto.createHash('sha256').update(key).digest('hex') });
  if (!device || device.status !== 'active') throw unauthorized('Unknown or inactive device', 'INVALID_DEVICE');

  const reading = await saveReading(device.patientId, req.body, 'device', device._id);
  device.lastSeenAt = new Date();
  await device.save();
  res.status(201).json({ success: true, data: { readingId: reading._id, abnormal: reading.abnormal } });
}));
\end{lstlisting}

\begin{howto}
A device is not a user: it has no password, no OTP and no JWT. It proves its
identity with its key, and the key's hash identifies the device and therefore
the patient. A device cannot choose which patient a reading belongs to, so a
stolen key can only add readings for one patient, and the admin can deactivate
the device at once. The Joi schema rejects physiologically impossible values
(for example a heart rate of 400).
\end{howto}

\subsection{The doctor's live monitor}

\code{GET \codesl{}api\codesl{}vitals\codesl{}monitor} returns, for every patient who currently allows
this doctor in, the latest reading and the patient's devices with their
\code{lastSeenAt}. The doctor's \emph{IoT Monitoring} page polls it, so readings
from devices (or from the simulator) appear within seconds.

\subsection{The simulator}

\file{scripts\codesl{}iot-simulator\codedot{}js} stands in for real hardware. It
reads the seeded device keys from \file{scripts\codesl{}\codedot{}demo-devices\codedot{}json} (git-ignored)
and posts readings to the real ingestion endpoint:

\begin{lstlisting}[language=JavaScript,caption={Simulated readings (\file{scripts/iot-simulator.js})}]
function readingFor(device) {
  if (spike) {
    return device.type === 'pulse_oximeter'
      ? { heartRate: Math.round(vary(128, 4)), spo2: Math.round(vary(89, 1)) }
      : { heartRate: Math.round(vary(124, 4)), systolic: Math.round(vary(168, 4)), diastolic: Math.round(vary(104, 3)), spo2: 96, temperature: vary(37.2, 0.2) };
  }
  if (device.type === 'pulse_oximeter') return { heartRate: Math.round(vary(80, 6)), spo2: Math.round(vary(97, 1)) };
  return {
    heartRate: Math.round(vary(74, 6)), systolic: Math.round(vary(126, 5)), diastolic: Math.round(vary(81, 3)),
    spo2: Math.round(vary(97.5, 1)), temperature: vary(36.8, 0.2),
  };
}
\end{lstlisting}

With \code{--spike}, the first reading from each device is clearly abnormal
(blood pressure around 168/104, SpO\textsubscript{2} around 89\%) and later ones
are normal again: a single burst that demonstrates the alert path end to end.
